The graph used for the transcriptomic analyses is distinguished from a subsequently enriched resource snapshot. Drug metadata were recovered from the OptimusKG drug parquet file and matched to Drug nodes by their source identifiers. Recovered fields included molecular structure, molecule type, maximum clinical phase, approval and withdrawal flags, and chemical descriptors. Where available, ChEMBL API records filled missing names, molecule types, phases or structures; OptimusKG values took precedence, with field-level source indicators preserving the distinction. Recovered names were stored separately from the original name property. Drug--gene evidence was recovered from the OptimusKG drug--gene parquet file and matched to original \texttt{ACTS\_ON} edges by drug and Ensembl gene identifiers. Available direct and indirect sources, relation assertions, conflict flags, reference URLs and mechanism text were retained. Derived evidence labels distinguished conflicting, multi-assertion, single-assertion and unattributed records; these were provenance summaries rather than independent experimental validation.

A later drug--gene extension was represented separately by \texttt{EXTERNAL\_ACTS\_ON} relationships. Newly introduced drug nodes carried both \texttt{Drug} and \texttt{ExternalDrug} labels; the latter marked a subset of the Drug category. External relationships retained source, release, action, mechanism and evidence fields. Their recorded provenance named ChEMBL, DGIdb, GtoPdb and Open Targets. Supplementary Table~\ref{tab:S4} distinguishes the analysis and enriched snapshots and reports populated-property coverage. These additions were not used to refit transcriptomic associations, clustering or candidate-gene scores. Snapshot verification compared all 8,044 progression relationships and their reported association properties with the final analysis export, and compared the 154 displayed gene--annotation connections by endpoint identifiers and relationship type.

